% Translated from sections/02-related-work.tex on 2026-09-23. The English file
% is canonical. Change it first, then carry the change here.
\section{Travaux connexes}
\label{sec:related}

\textbf{Reconnaissance.} CATMuS est une famille de jeux de données et de
modèles dont les transcriptions suivent toutes les mêmes règles, si bien que
les modèles entraînés dessus produisent des sorties comparables. Elle est née
avec les manuscrits médiévaux~\cite{clerice2024catmusmedieval}, et CATMuS-Print
applique la même démarche aux livres imprimés~\cite{gabay2024catmusprint}. J'ai
affiné le modèle CATMuS-Print [Large]~\cite{gabay2024catmusprintlarge} avec
kraken~7.0.2. L'article de kraken décrit la version~5~\cite{kiessling2025kraken}
et la notice Zenodo fixe la version 3.0b20~\cite{kiessling2021krakensoftware};
la version utilisée est donc donnée par son numéro. Les pages d'entraînement
ont été transcrites dans Transkribus~\cite{kahle2017transkribus,
muehlberger2019transkribus}. Les travaux publiés les plus proches sur un
matériau semblable sont CATMuS Gothic Print, un autre modèle de base entraîné
sur des textes religieux et polémiques~\cite{solfrini2024catmusgothic}, et
l'évaluation par Tiger de l'affinage de modèles HTR sur le français des
\textsc{xvii}\textsuperscript{e} et \textsc{xviii}\textsuperscript{e}~siècles~\cite{tiger2026overfit}.

\textbf{Correction par modèles de langue.} L'étape de correction a été conçue
autour d'une mise en garde qui revient sans cesse dans la littérature. Kanerva
et al.\ montrent qu'un correcteur à base de LLM peut dégrader une bonne
reconnaissance~\cite{kanerva2025nofreelunches}, et Boros et al.\ constatent la
même chose sur des documents historiques, avec plusieurs manières de formuler
la requête~\cite{boros2024postcorrection}. Cela compte ici parce que le modèle
de reconnaissance est déjà précis: avec un CER de \trainCer{} sur ses pages de
validation, presque chaque caractère qu'il produit est juste, et un modèle qui
relit des lignes entières a plus d'occasions d'abîmer le texte que de le
réparer. Ma passe n'envoie donc que les mots dont le modèle n'était pas sûr, et
ne garde une réponse que si le reste de la ligne revient inchangé. Sur ce
matériau, la surcorrection typique est la modernisation: un modèle à qui l'on
demande de corriger un \qtr{ſ} mal lu met souvent aussi à jour l'orthographe
autour, et la règle d'acceptation refuse ces réponses. La compétition
HIPE-OCRepair décrit la surcorrection d'un texte propre comme son problème le
plus persistant, et note que les taux d'erreur caractère ne la montrent pas
entièrement~\cite{ehrmann2026hipeocrepair}. C'est pourquoi la
section~\ref{sec:e2} rapporte ce que la règle refuse autant que ce qu'elle
garde.

D'autres travaux placent après la reconnaissance un modèle de correction
entraîné pour cette tâche. Un module ByT5 placé après kraken réduit le CER sur
des incunables du \textsc{xv}\textsuperscript{e}~siècle~\cite{momtaz2025modular},
et Thomas et al.\ étudient la correction par LLM de journaux
anciens~\cite{thomas2024leveraging}. Ma passe utilise au contraire un modèle
ouvert généraliste, gpt-oss 20B, exécuté en local sans entraînement pour la
tâche, et s'en remet à la règle d'acceptation pour le tenir.

Ces études mesurent la correction comme une variation du taux d'erreur. Cela
décrit la sortie du correcteur, mais pas la règle qui décide de la garder. Je
n'ai trouvé aucun travail publié qui rapporte ce qu'une telle règle refuse,
même si une recherche bibliographique ne peut exclure une thèse ou une
publication non anglophone qui m'aurait échappé. Les intervalles de confiance
de ce rapport suivent la méthode du bootstrap par
percentiles~\cite{efron1994bootstrap}.
